% ATTEMPT_STATUS: unresolved
% RESEARCH_GENERATED: 2026-10-01; GPT-6 Astra Ultra; individual mathematical attempt
% TOP_PROBLEM: {"id": 111, "title": "Large Cosets in Iterated Sumsets", "category_id": 2, "category": {"id": 2, "name": "combinatorics", "display_name": "Combinatorics", "description": "Counting problems, graph theory, discrete structures.", "slug": "combinatorics", "order_index": 2, "created_at": "2026-07-31T15:26:25.671Z"}, "status": "open"}
% FIRST_POSTED: 2026-10-01
% Imported from: unsolvedmath_additional_500.tex; UnsolvedMath ID 111
% !TeX program = lualatex
% Compile twice: lualatex 111.tex
% Self-contained: no external bibliography, images, or \input files.
% Numeric section labels and visible numbers are original UnsolvedMath IDs.
\documentclass[11pt,a4paper]{article}
\usepackage[margin=24mm,headheight=14pt]{geometry}
\usepackage{fontspec}
\setmainfont{Latin Modern Roman}
\setsansfont{Latin Modern Sans}
\setmonofont{Latin Modern Mono}
\usepackage{amsmath,amssymb,mathtools}
\usepackage{unicode-math}
\setmathfont{Latin Modern Math}
\usepackage{microtype}
\usepackage{enumitem,xurl,needspace}
\usepackage[dvipsnames]{xcolor}
\usepackage{hyperref}
\hypersetup{unicode=true,colorlinks=true,linkcolor=MidnightBlue,urlcolor=MidnightBlue,
 pdftitle={UnsolvedMath Problem 111},pdfauthor={Source-annotated compilation},bookmarksnumbered=true}
\usepackage{bookmark,tocloft,titlesec,fancyhdr}
\setlength{\cftsecnumwidth}{4.2em}
\cftsetpnumwidth{2.5em}
\cftsetrmarg{4em}
\titleformat{\section}{\Large\bfseries\raggedright}{\thesection}{0.7em}{}
\pagestyle{fancy}\fancyhf{}
\fancyhead[L]{\small\sffamily UnsolvedMath research attempt}
\fancyhead[R]{\small Original catalogue IDs}
\fancyfoot[C]{\thepage}
\setlength{\parindent}{0pt}
\setlength{\parskip}{5pt plus 1pt}
\setlength{\emergencystretch}{3em}
\setcounter{tocdepth}{1}\setcounter{secnumdepth}{1}
\setlist[itemize]{leftmargin=3em,itemsep=4pt,topsep=4pt,parsep=0pt}
\allowdisplaybreaks
\newcommand{\N}{\mathbb{N}}
\newcommand{\Z}{\mathbb{Z}}
\newcommand{\Q}{\mathbb{Q}}
\newcommand{\R}{\mathbb{R}}
\newcommand{\C}{\mathbb{C}}
\newcommand{\F}{\mathbb{F}}
\newcommand{\A}{\mathbb{A}}
\newcommand{\PP}{\mathbb{P}}
\newcommand{\E}{\mathbb{E}}
\newcommand{\eps}{\varepsilon}
\newcommand{\dd}{\,\mathrm d}
\newcommand{\sref}[2]{\hyperlink{source:#1:#2}{[S#2]}}
\newif\ifshowcatalogue
\showcataloguetrue
\newcommand{\cataloguescope}[1]{\ifshowcatalogue
 \paragraph{Statement recorded in the supplied source.}
 #1\fi}
\begin{document}
% UnsolvedMath ID 111; source catalogue code GREEN-023

\setcounter{section}{110}

\section{Large Affine Subspaces in a Tenfold Sumset}\label{111}

{\small\sffamily UnsolvedMath ID 111 · Combinatorics}

\subsection{Definitions and mathematical statement}

\paragraph{Shared notation.}Unless a section says otherwise, $\N=\{1,2,\ldots\}$,
$\Z$ is the ring of integers, $\Q,\R,\C$ are the rational, real, and complex
fields, $[N]=\{1,\ldots,\lfloor N\rfloor\}$, and $\log$ is the natural
logarithm. For real functions of a parameter tending to infinity,
$f\ll g$ and $f=O(g)$ mean $|f|\leq Cg$ eventually for some fixed $C$;
$f\gg g$ reverses this comparison; $f=o(g)$ means $f/g\to0$;
$f\sim g$ means $f/g\to1$; and $f\asymp g$ means both order bounds.
A subscript on $O$ or $\ll$ specifies allowed constant dependence.
The expression $n^{o(1)}$ in an upper bound means growth smaller than
$n^\epsilon$ for each fixed $\epsilon>0$. Local definitions govern any
exception, finite-field convention, or use of the same symbol for another
object. An asymptotic claim is not automatically an effective algorithm.



Here $10A=\{a_1+\cdots+a_{10}:a_i\in A\}$, with repetitions allowed. A coset is an affine subspace $x+V$, where $V\leq\F_2^n$ is linear, and its codimension is $n-\dim V$. The bound asks for an absolute $C$ such that one can achieve $\operatorname{codim}V\leq C\log(1/\alpha)$ for every nonempty $A$ of density $\alpha$. \sref{111}{1} \sref{111}{2}

\paragraph{Statement recorded in the supplied source.}
Suppose that $A \subset \mathbb{F}_2^n$ has density $\alpha$. Does $10A$ contain a coset of some subspace of dimension at least $n - O(\log(1/\alpha))$? \sref{111}{1}

\paragraph{Residual task reported by the archive.}

The embedded review (2026-08-17) records: Show that a fixed bounded iterated sumset, specifically 10A, contains a coset of codimension O(log(1/alpha)). \sref{111}{1}

\subsection{Short English statement}

Must the sum of ten copies of a dense set over the two-element field contain an affine subspace with only logarithmic codimension loss?

\subsection{Research attempt}

\paragraph{Provenance and scope.}
This individual attempt was generated by GPT-6 Astra Ultra on 1 October 2026. The method was to derive explicit partial results and test the first proposed extension adversarially. The original statement and sources below are retained. No claim of mathematical novelty or complete solution is made.

\paragraph{Formal target and current scope.}
Let $V=\F_2^n$ and $A\subset V$ have density $\alpha>0$. The notation $10A$ allows repetitions. The question asks for a coset of codimension $O(\log(1/\alpha))$, with an absolute implied constant. The primary problem list explicitly separates this polynomial Bogolyubov statement from the solved polynomial Freiman--Ruzsa conjecture. The inspected Marton-conjecture paper is consequently not being cited as a solution of this problem. We reconstruct a quantitative Fourier argument and isolate the exact loss preventing the desired logarithmic codimension.

\paragraph{Normalisation and the spectral cutoff.}
For $\xi\in V$ write $\chi_\xi(x)=(-1)^{\xi\cdot x}$ and $\widehat f(\xi)=\E_xf(x)\chi_\xi(x)$, with $f=1_A$. For normalised convolution, Fourier coefficients multiply. Parseval gives $\sum_\xi\widehat f(\xi)^2=\alpha$, and $\widehat f(0)=\alpha$. Let
\[
 \tau^2=\alpha^3/2,\qquad
 \Gamma=\{\xi:|\widehat f(\xi)|\geq\tau\},\qquad
 H=\{x:\xi\cdot x=0\text{ for every }\xi\in\Gamma\}.
\]
There are at most $\alpha/\tau^2=2\alpha^{-2}$ frequencies in $\Gamma$, so $\operatorname{codim}H\leq2\alpha^{-2}$. This dimension estimate uses the number of equations as an upper bound on their rank; dependencies could improve it but are not assumed.

For $x\in H$, Fourier inversion gives
\[
 f^{*4}(x)=\sum_\xi\widehat f(\xi)^4\chi_\xi(x)
 \geq\alpha^4-\sum_{\xi\notin\Gamma}\widehat f(\xi)^4
 \geq\alpha^4-\tau^2\sum_\xi\widehat f(\xi)^2
 =\alpha^4/2>0.
\]
The zero frequency belongs to $\Gamma$ because $\alpha^2\geq\alpha^3/2$, and every retained character equals one on $H$. Positivity of the convolution means that there are four elements of $A$ summing to $x$. Therefore $H\subset4A$. Choose any fixed $a\in A$ and append three pairs $(a,a)$ to any such representation. Since $a+a=0$, this proves $4A\subset10A$. We have obtained a subspace, not merely a coset, of codimension at most $2\alpha^{-2}$.

\paragraph{A density-amplification lemma and its exact threshold.}
If $B\subset t+W$ occupies more than half the coset, then $B+B$ contains all of $W$. For $w\in W$, both $B$ and $w+B$ lie in $t+W$ and their cardinalities sum to more than $|W|$, so they intersect. An intersection point gives $w=b+b'$ with $b,b'\in B$. This proof also shows that density strictly greater than one half is necessary for this elementary conclusion: a subspace of index two inside $W$ has density one half but its double sumset is itself.

Applied to $B=2A\cap(t+W)$, the lemma would give $W\subset4A$ once a suitable high-density coset has been found. Thus a sufficient route to the requested bound is to find a codimension $O(\log(1/\alpha))$ coset on which $2A$ has density greater than one half. The lemma proves the implication; existence of that coset is not supplied by the preceding Fourier cutoff.

\paragraph{The exponent loss cannot be erased by notation.}
The cutoff proof controls the spectral tail by $\tau^2\alpha$. To make this smaller than the main term $\alpha^4$, it requires $\tau^2<\alpha^3$. Parseval then only bounds the number of selected frequencies at scale $\alpha^{-2}$. Replacing this cardinality by $O(\log(1/\alpha))$ without a new rank theorem is unsupported. A large spectrum can have many additive dependencies, but proving enough of them at the needed threshold is precisely additional mathematical content.

There is also a matching lower obstruction to improving the target below logarithmic order: if $A$ is a subspace of codimension $k$, then $\alpha=2^{-k}$ and $10A=A$. Every affine subspace contained in $A$ has dimension at most $n-k$. Thus at least $\log_2(1/\alpha)$ codimension is sometimes unavoidable. This checks that the catalogue's proposed scale has the right basic lower benchmark.

\paragraph{Verification and unresolved step.}
For $\alpha=1$ the spectral construction returns the whole space. For a singleton it may give a vacuous numerical dimension bound, but the convolution and containment remain valid. Repeated summands are essential to the passage from four to ten copies and are expressly allowed by the statement. The proof uses no unverified consequences of the Marton theorem. The established result is the explicit $2\alpha^{-2}$ codimension bound together with the density-amplification criterion. The missing step is a logarithmic-rank structural improvement or a comparably strong dense-coset argument; the full target remains unresolved here.

\subsection{Sources}

{\small

\begin{itemize}

\item[\hypertarget{source:111:1}{S1}] ULAM.AI, \emph{UnsolvedMath}, supplied \texttt{problems.json}, record 111 (GREEN-023), statement and background. The embedded literature-review date is 2026-08-17; that is the archive’s review date, not a fresh certification. Dataset landing page: \url{https://huggingface.co/datasets/ulamai/UnsolvedMath}.

\item[\hypertarget{source:111:2}{S2}] Tom Sanders, On the Bogolyubov-Ruzsa lemma, arXiv:1011.0107; Analysis \& PDE 5 (2012), 627-655. \url{https://arxiv.org/abs/1011.0107}. Bibliographic information imported from the supplied record.

\item[\hypertarget{source:111:3}{S3}] Tomasz Kościuszko, Counting solutions to invariant equations in dense sets, arXiv:2306.08567, Theorem 9. \url{https://arxiv.org/abs/2306.08567}. Bibliographic information imported from the supplied record.

\item[\hypertarget{source:111:4}{S4}] Ben Green, 100 Open Problems, Problem 50 and 2024 update, . \url{https://people.maths.ox.ac.uk/greenbj/papers/open-problems.pdf}. The author-hosted document was inspected on 27 September 2026; the archive’s local code is not a problem-number locator in this paper.

\item[E1] W. T. Gowers, B. Green, F. Manners and T. Tao, \emph{On a conjecture of Marton}. \url{https://arxiv.org/abs/2311.05762}. Inspected 1 October 2026.

\end{itemize}

}

\label{end:111}
\end{document}
